
\section{特别是概形上的\trLocallyFreeSheaf}

我们现在来讨论\trLocallyFreeSheaf{}的若干更专门针对\trScheme{}的面向。出于你对\trManifold{}上\trLineBundle{}的直觉，你或许会希望：每个点都有一个「小」的开\trNeighborhood，使得在其上一切\trInvertibleSheaf（或\trLocallyFreeSheaf）都平凡。令人遗憾的是，事实并非如此。我们终将在（\S19.11.1）看到：对于 $\mathbb{A}^2_{\mathbb{C}}$ 中由 $y^2-x^3-x=0$ 给出的这条\trCurve，每个非空\trOpenSet{}上都有非平凡的\trInvertibleSheaf。（这将用到「它是某条\emph{\trEllipticCurve}的开子集」这一事实。）这是\trAlgebraicGeometry{}中一个挥之不去的问题：\trZariskiTopology{}太钝了。

\textbf{14.2.A.（容易的练习）}\quad 证明：\trScheme{}上的\trLocallyFreeSheaf{}都是\trQC{}的。

类似地，有限\trRank{}的\trLocallyFreeSheaf{}总是\trFiniteType{}的；而若 $\OO_X$ 是\trCoh{}的，则 $X$ 上有限\trRank{}的\trLocallyFreeSheaf{}也是\trCoh{}的。（记住：若 $\OO_X$ 不是\trCoh{}的，那么凝聚性在 $X$ 上就是个相当无用的观念。）

\textbf{14.2.B.（较不重要的练习）}\quad 用\mbox{练习~6.1.B(b)}中的例子证明：并非每个\trQCSheaf{}都是\trLocallyFreeSheaf。

\textbf{14.2.C.（练习）}\quad 证明：$\mathbb{A}^1_k$ 上每个（有限）\trRank{}为 $n$ 的\trVectorBundle{}都是\trRank{}为 $n$ 的\trTrivialBundle。提示：主理想整环上的有限生成模是\trCyclicModule{}的有限直和，如注记~6.4.4 所述。（当心：在援引这条结果之前，你得先确信所论的\trModule{}确实有限生成！）向 $\mathbb{A}^n_k$ 的困难推广，见 \S15.4.12 中的旁注。

\textbf{14.2.1.}\ \emph{\trhDefinition{}。}\quad \trScheme{} $X$ 上\trLocallyFreeSheaf{}的\textbf{有理（及正则）\trSection}，其定义一如我们当年对有理（及正则）函数的定义（例如见 \S6.6 与 \S7.5）。

\textbf{14.2.D.（练习）}\quad 设 $s$ 是\trScheme{} $X$ 上\trLocallyFreeSheaf{} $\scr{F}$ 的一个\trSection。定义\textbf{由 $s=0$ 切出的\trClosedSubscheme}这一观念，记作 $V(s)$（理由显然）。务必检验你的定义不依赖于种种选取！（这道练习为\cref{ex:vanishing-scheme-definition}给出一种新解法。）提示：取某个开集 $U$ 上的一个\trTrivialization，$s$ 便对应于 $U$ 上的一组函数 $f_1,\dots$；在 $U$ 上，取由这些函数切出的概形。另一条避开坐标的提示：想办法把它定义为使 $s$ 限制为 $0$ 的最大\trClosedSubscheme。

\textbf{14.2.E.（容易的练习）}\quad 证明：局部\trNoetherian{}正规概形上的\trLocallyFreeSheaf{}满足「代数\trHartogs{}引理」：在一个\trCodimension{}至少为 $2$ 的集合之外有定义的\trSection，都可以越过该集合延拓。（\trNoetherian{}正规概形的代数\trHartogs{}引理即定理~13.5.19。）

\textbf{14.2.F.（容易的练习）}\quad 设 $s$ 是局部\trNoetherian{}正规概形上某个\trInvertibleSheaf{}的一个非零有理\trSection。证明：若 $s$ 无\trPole，则 $s$ 是正则的。（提示：\mbox{练习~13.5.I}。）

\textbf{14.2.2.\ 短正合列中的\trLocallyFreeSheaf。}

\textbf{14.2.G.（练习：对日后的问题或许有帮助）}\quad
\begin{enumerate}[label=(\alph*)]
\item 设
\begin{equation*}\tag{14.2.2.1}\label{eq:qcoh-short-exact-sequence}
0\longrightarrow\scr{F}'\longrightarrow\scr{F}\longrightarrow\scr{F}''\longrightarrow 0
\end{equation*}
是 $X$ 上\trQCSheaf{}的一个\trShortExactSequence。设 $U=\Spec A$ 是一个使 $\scr{F}'$ 与 $\scr{F}''$ 自由的仿射开集，比如说 $\scr{F}'|_{\Spec A}=\widetilde{A}^{\oplus a}$、$\scr{F}''|_{\Spec A}=\widetilde{A}^{\oplus b}$。（这里为方便起见假定 $a$ 与 $b$ 有限，但这并非必要，所以尽管放心推广到无限\trRank{}的情形。）证明：$\scr{F}$ 在 $\Spec A$ 上也自由，并且 $0\to\scr{F}'\to\scr{F}\to\scr{F}''\to 0$ 可以解读为来自重言正合列 $0\to A^{\oplus a}\to A^{\oplus(a+b)}\to A^{\oplus b}\to 0$。（作为推论：给定\trQCSheaf{}的一个正合列 \eqref{eq:qcoh-short-exact-sequence}，若其中 $\scr{F}'$ 与 $\scr{F}''$ 都\trLocallyFree，则 $\scr{F}$ 也必定\trLocallyFree。）
\item 在有限\trRank{}的情形，证明：给定一个由使之\trTrivialization{}的仿射开集（即 (a) 中所述形状的开集）组成的\trOpenCover，$\scr{F}$ 的\trTransitionFunction（其实是矩阵）可以解读为块上三角矩阵，其中左上角的诸 $a\times a$ 块是 $\scr{F}'$ 的\trTransitionFunction，而底部的诸 $b\times b$ 块是 $\scr{F}''$ 的\trTransitionFunction。
\end{enumerate}

\textbf{14.2.H.（练习）}\quad 给定 $X$ 上\trQCSheaf{}的\trShortExactSequence \eqref{eq:qcoh-short-exact-sequence}。由\mbox{练习~14.2.G(a)}，若 $\scr{F}'$ 与 $\scr{F}''$ 都\trLocallyFree，则 $\scr{F}$ 也\trLocallyFree。
\begin{enumerate}[label=(\alph*)]
\item 若 $\scr{F}$ 与 $\scr{F}''$ 是\emph{有限\trRank}的\trLocallyFreeSheaf，证明 $\scr{F}'$ 也是。提示：化归到 $X=\Spec A$ 而 $\scr{F}$ 与 $\scr{F}''$ 自由的情形。把映射 $q\colon\scr{F}\to\scr{F}''$ 解读为一个在 $A$ 中取值的 $n\times m$ 矩阵 $M$，其中 $m$ 是 $\scr{F}$ 的\trRank，$n$ 是 $\scr{F}''$ 的\trRank。对 $X$ 的每一点 $p$，证明：$M$ 存在 $n$ 个列 $\{c_1,\dots,c_n\}$，它们在 $p$ 处线性无关，从而在 $p$ 附近也线性无关（因为线性无关由相应 $n\times n$ \trDeterminant{}的非零性给出）。于是，$X$ 可以被一些\trDistinguishedOpen{}覆盖，这些开集与「$M$ 的 $n$ 列的诸选取」一一对应。限制到其中一个子集并给诸列改名，化归到「$M$ 前 $n$ 列的\trDeterminant{}可逆」的情形。然后在 $A^{\oplus m}=\scr{F}(\Spec A)$ 上换\trCoordinate，使得 $M$ 在新\trCoordinate{}下前 $n$ 列为单位矩阵、其后为 $0$。最后，在这种情形下把 $\scr{F}'$ 解读为 $\widetilde{A^{\oplus(m-n)}}$。
\item 若 $\scr{F}'$ 与 $\scr{F}$ 都\trLocallyFree，证明 $\scr{F}''$ 未必\trLocallyFree。提示：考虑 \eqref{eq:tk-k-modules-exact}，我们可以把它解读为 $\mathbb{A}^1$ 上某一点的闭子概形正合列 \eqref{eq:closed-subscheme-exact-sequence}。（若\trLocallyFreeSheaf{}之间的一个单射 $\alpha\colon\scr{F}'\hookrightarrow\scr{F}$ 使其余核 $\scr{F}''$ 也\trLocallyFree，人们有时就说 $\alpha$ 是一个\emph{\trVectorBundle{}的单射}（injection of vector bundles），而不只是\trLocallyFreeSheaf{}的单射。他们暗中考虑的是一个有别于\trLocallyFreeSheaf{}范畴的\trVectorBundle{}范畴；在这样的范畴里，\trVectorBundle{}之间映射的余核总还是\trVectorBundle。进一步的讨论见 \S14.2.4。）
\end{enumerate}

\textbf{14.2.I.（练习）}\quad
\begin{enumerate}[label=(\alph*)]
\item 设 $0\to\scr{F}'\to\scr{F}\to\scr{F}''\to 0$ 是有限\trRank{}的\trLocallyFreeSheaf{}的一个\trShortExactSequence。证明：
\[
0\longrightarrow\scr{F}''^{\vee}\longrightarrow\scr{F}^{\vee}\longrightarrow\scr{F}'^{\vee}\longrightarrow 0
\]
（由对原正合列施加 $\scr{H}om(\cdot,\OO_X)$ 而得）也是有限\trRank{}的\trLocallyFreeSheaf{}的一个\trShortExactSequence。
\item[\hard~(b)] 设 $\cdots\to\scr{F}_2\to\scr{F}_1\to\scr{F}_0\to 0$ 是有限\trRank{}的\trLocallyFreeSheaf{}的一个正合列。证明：由此诱导的列
\[
0\longrightarrow\scr{F}_0^{\vee}\longrightarrow\scr{F}_1^{\vee}\longrightarrow\scr{F}_2^{\vee}\longrightarrow\cdots
\]
也是有限\trRank{}的\trLocallyFreeSheaf{}的一个正合列。
\end{enumerate}

\textbf{14.2.3.\ \trVectorBundle{}的映射。}

在更通常的意义下，\trVectorBundle{}之间的映射比那些碰巧是\trLocallyFreeSheaf{}的\trQCSheaf{}之间的\trMorphism{}限制更强。一个\trLocallyFreeSheaf{}的\trLocallyFree\trSubsheaf，并不总能给出某个\trVectorBundle{}的一个「子向量丛」（subvector bundle）。最典型的例子是 $\mathbb{A}^1_k=\Spec k[t]$ 上与如下 $k[t]$-模正合列相对应的那个\trQCSheaf{}正合列：
\begin{equation*}\tag{14.2.3.1}\label{eq:tk-k-modules-exact}
0\longrightarrow tk[t]\longrightarrow k[t]\longrightarrow k\longrightarrow 0.
\end{equation*}
\trLocallyFreeSheaf $\widetilde{tk[t]}$ 是 $\widetilde{k[t]}$ 的一个\trSubsheaf，但它并不对应于一个「子向量丛」；其余核不是一个\trVectorBundle。（这个例子曾在 \S14.1.8 中提及。）

由于真正「是\trVectorBundle{}之间的映射」的\trLocallyFreeSheaf{}映射有着某些格外良好而有用的性态，似乎值得给它们起个名字。这种映射并没有公认的名字，所以我们就自己挑一个。

\textbf{14.2.4.}\ \emph{\trhDefinition{}。}\quad 设 $X$ 是一个\trScheme（或者，其实可以更一般地，是一个\trRingedSpace），并设 $\scr{E}$ 与 $\scr{F}$ 是 $X$ 上两个有限\trRank{}的\trLocallyFreeSheaf，\trRank{}分别为 $a+b$ 与 $a+c$。我们说 $\phi\colon\scr{E}\to\scr{F}$ 是一个\textbf{秩 $a$ 的\trVectorBundle{}映射}（rank $a$ map of vector bundles），若对每一点 $p\in X$，都存在 $p$ 的一个\trOpen\trNeighborhood{} $U$，使之带有交换图
\[
\begin{tikzcd}[row sep=3em]
\scr{E}|_U \ar[d, "\sim"'] \ar[rr, "\phi"] & & \scr{F}|_U \ar[d, "\sim"'] \\
\OO_U^{\oplus(a+b)} \ar[r, two heads] & \OO_U^{\oplus a} \ar[r, hook] & {\OO_U^{\oplus(a+c)},}
\end{tikzcd}
\]
其中映射 $\OO_U^{\oplus(a+b)}\twoheadrightarrow\OO_U^{\oplus a}$ 是到前 $a$ 个直和项上的\trProjection，而映射 $\OO_U^{\oplus a}\hookrightarrow\OO_U^{\oplus(a+c)}$ 是作为前 $a$ 个直和项的\trInclusionMap。若在每一点 $p\in X$ 附近，$\phi$ 都是秩 $a$ 的\trVectorBundle{}映射（对某个 $a$），我们就说 $\phi$ 是一个\textbf{\trVectorBundle{}的映射}（map of vector bundles）。若 $b=0$，我们就说 $\phi$ 是一个\textbf{\trVectorBundle{}的单射}（如\mbox{练习~14.2.H(b)}中那样）；我们对这一称谓的使用将很节制。

\textbf{14.2.5.}\ \emph{重要但直截了当的观察。}\quad (i) 「秩 $a$ 的\trVectorBundle{}映射」这一观念与任何\trBaseChange{} $Y\to X$ 相交换。(ii) 诸\trQCSheaf{} $\ker\phi$、$\im\phi$ 与 $\coker\phi$ 都\trLocallyFree（其\trRank{}分别为 $b$、$a$ 与 $c$，皆有限），而且它们的构造与任何\trBaseChange{} $Y\to X$ 相交换。
\[
\begin{tikzcd}[row sep=3em, column sep=2.4em]
0 \ar[r] & \ker\phi \ar[r] \ar[d, "\sim"'] & \scr{E} \ar[r, "\phi"] \ar[d, "\sim"'] & \scr{F} \ar[r] \ar[d, "\sim"'] & \coker\phi \ar[r] \ar[d, "\sim"'] & 0 \\
0 \ar[r] & \OO_U^{\oplus b} \ar[r] & \OO_U^{\oplus(a+b)} \ar[r] & \OO_U^{\oplus(a+c)} \ar[r] & \OO_U^{\oplus c} \ar[r] & 0
\end{tikzcd}
\]
(iii) 若 $\phi\colon\scr{E}\to\scr{F}$ 是秩 $a$ 的\trVectorBundle{}映射，则在任意一点 $p\in X$ 处，$\phi|_p\colon\scr{E}|_p\to\scr{F}|_p$ 的\trRank{}为 $a$。

在接下来的几道练习中，我们给出三条判别法，用来判定有限\trRank{}的\trLocallyFreeSheaf{}之间的一个\trMorphism{} $\phi$ 何时是\trVectorBundle{}的映射。

\textbf{14.2.J.（练习：「\trVectorBundle{}的映射」的一个有用刻画）}\quad 设 $\phi\colon\scr{E}\to\scr{F}$ 是有限\trRank{}的\trLocallyFreeSheaf{}之间的一个\trMorphism，其中 $\scr{F}$ 的\trRank{}为 $a+c$。证明：$\phi$ 是秩 $a$ 的\trVectorBundle{}映射，当且仅当 $\coker\phi$ 是\trRank{}为 $c$ 的\trLocallyFreeSheaf。（特别地，有限\trRank{}的\trLocallyFreeSheaf{}之间的每个\trSurjectiveMap{}都是\trVectorBundle{}的映射。）

\textbf{14.2.K.（练习）}\quad 设 $\phi\colon\scr{E}\to\scr{F}$ 是有限\trRank{}的\trLocallyFreeSheaf{}之间的一个\trMorphism，并设 $p\in X$。证明：自然映射 $(\ker\phi)|_p\to\ker(\phi|_p)$ 是\trSurjectiveMap，当且仅当 $\phi$ 在 $p$ 的某个\trNeighborhood{}内是\trVectorBundle{}的映射。

\textbf{14.2.6.}\ \hard\ \textbf{张量\trAlgebra{}构造。}

（这些重要的想法需要更多的工夫，所以你不妨把它们留待日后。）为接下来的几道练习，我们回忆如下内容。若 $M$ 是一个 $A$-\trModule，则\textbf{张量\trAlgebra}（tensor algebra）$\mathbf{T}^\bullet(M)$ 是一个非\trCommutative{}代数，以 $\mathbb{Z}^{\ge 0}$ \trGrading，定义如下：$\mathbf{T}^0(M)=A$，$\mathbf{T}^n(M)=M\otimes_A\cdots\otimes_A M$（乘积中出现 $n$ 个项），而乘法一如你的预期。

\trSymmetricAlgebra{} $\Sym^\bullet M$ 是一个\trSymmetricAlgebra，以 $\mathbb{Z}^{\ge 0}$ \trGrading，定义为 $\mathbf{T}^\bullet(M)$ 对由一切形如 $x\otimes y-y\otimes x$（其中 $x$、$y\in M$）的元素所生成的（双边）\trIdeal{}的商。于是，$\Sym^n M$ 是 $M\otimes\cdots\otimes M$ 对形状如 $m_1\otimes\cdots\otimes m_n-m_1'\otimes\cdots\otimes m_n'$ 的诸\trRelation{}的商，其中 $(m_1',\dots,m_n')$ 是 $(m_1,\dots,m_n)$ 的一个重排。

\textbf{\trExteriorAlgebra}（exterior algebra）$\wedge^\bullet M$ 定义为 $\mathbf{T}^\bullet M$ 对由一切形状为 $x\otimes x$（$x\in M$）的元素所生成的（双边）\trIdeal{}的商。展开 $(a+b)\otimes(a+b)$，我们便看到在 $\wedge^2 M$ 中 $a\otimes b=-b\otimes a$。这蕴涵：若 $2$ 在 $A$ 中\trInvertible（例如若 $A$ 是一个\trCharacteristic{}不为 $2$ 的\trField），则 $\wedge^n M$ 是 $M\otimes\cdots\otimes M$ 对形状如 $m_1\otimes\cdots\otimes m_n-(-1)^{\operatorname{sgn}(\sigma)}m_{\sigma(1)}\otimes\cdots\otimes m_{\sigma(n)}$ 的诸\trRelation{}的商，其中 $\sigma$ 是 $\{1,\dots,n\}$ 的一个置换。\trExteriorAlgebra{}是一个「反交换」的 $A$-\trAlgebra。

更好的一点是：$\Sym$ 与 $\wedge$ 都可以由\trUniversalProperty{}来定义。例如，任何 $A$-模的映射 $M^{\otimes n}\to N$，只要在其 $n$ 个入口上是对称的，就唯一地穿过 $\Sym^n_A(M)$。

最正确的写法是 $\mathbf{T}^\bullet_A(M)$、$\Sym^\bullet_A(M)$ 与 $\wedge^\bullet_A(M)$，但出于方便，「基环」$A$ 通常被省去。

\textbf{14.2.L.（练习）}\quad 若 $\scr{F}$ 是\trRank{}为 $m$ 的\trLocallyFreeSheaf，证明 $\mathbf{T}^n\scr{F}$、$\Sym^n\scr{F}$ 与 $\wedge^n\scr{F}$ 都\trLocallyFree，并求出它们的\trRank。我们指出：在这种情形，$\wedge^{\operatorname{rank}\scr{F}}\scr{F}$ 被记作 $\det\scr{F}$，称为\textbf{\trDeterminant{}（线）丛}（determinant (line) bundle），或者（既更好又更糟）称为\textbf{\trDeterminant{}\trLocallyFreeSheaf}（determinant locally free sheaf）。

（评注：这些构造可以对任意\trRingedSpace{}上的 $\OO$-\trModule{}来定义。）

你还可以定义非\trCommutative{}\trAlgebra{}的层 $\mathbf{T}^\bullet\scr{F}$、\trCommutative{}\trAlgebra{}的层 $\Sym^\bullet\scr{F}$，以及「反交换」\trAlgebra{}的层 $\wedge^\bullet\scr{F}$。

\textbf{14.2.M.（练习）}\quad 设 $0\to\scr{F}'\to\scr{F}\to\scr{F}''\to 0$ 是\trLocallyFreeSheaf{}的一个正合列。证明：对任意 $r$，$\Sym^r\scr{F}$ 有一个\trFiltration
\[
\Sym^r\scr{F}=\scr{G}^0\supset\scr{G}^1\supset\cdots\supset\scr{G}^r\supset\scr{G}^{r+1}=0,
\]
其\trSubquotient{}为
\[
\scr{G}^p/\scr{G}^{p+1}\cong(\Sym^p\scr{F}')\otimes(\Sym^{r-p}\scr{F}'').
\]
（这里给出两条不同的可行提示，供本题与\mbox{练习~14.2.N}使用：(1) 把 $\scr{F}$ 的转移矩阵解读为分两块的块上三角矩阵，其中一个对角块给出 $\scr{F}'$ 的转移矩阵，另一个给出 $\scr{F}''$ 的转移矩阵（参看\mbox{练习~14.2.G(b)}。\footnote{原书此处作（cf.\ Exercise 14.2.2.1(b)），而 14.2.2.1 系公式编号、并无 (b) 小问；按上下文（转移矩阵分块上三角）显指练习 14.2.G(b)，径改。——译者注}）然后适当地把 $\Sym^r\scr{F}$ 的转移矩阵解读为分 $r+1$ 块的块上三角矩阵。(2) 只需考虑一个足够小的仿射开集 $\Spec A$（使 $\scr{F}'$、$\scr{F}$、$\scr{F}''$ 在其上自由），并证明你的构造对在任一元素 $f\in A$ 处取\trLocalization{}表现良好，这就够了。在这样的开集中，由\mbox{练习~14.2.G}，该正合列为 $0\to A^{\oplus p}\to A^{\oplus(p+q)}\to A^{\oplus q}\to 0$。设 $e_1,\dots,e_p$ 是 $A^{\oplus p}$ 的标准基，$f_1,\dots,f_q$ 是 $A^{\oplus q}$ 的标准基。设 $e_1',\dots,e_p'$ 表示 $e_1,\dots,e_p$ 在 $A^{\oplus(p+q)}$ 中的像。设 $f_1',\dots,f_q'$ 是 $f_1,\dots,f_q$ 到 $A^{\oplus(p+q)}$ 中的任一组提升。\footnote{原书本段三处显见排印笔误：「by the Exercise 14.2.G」衍「the」、「the the standard basis」叠「the」、「be denote the images」衍「be」，均径改。——译者注}注意 $f_i'$ 在模 $e_1',\dots,e_p'$ 之下是\trWellDefined{}的。注意
\[
\Sym^r\scr{F}|_{\Spec A}\cong\bigoplus_{i=0}^{r}\Sym^i\scr{F}'|_{\Spec A}\otimes_{\OO_{\Spec A}}\Sym^{r-i}\scr{F}''|_{\Spec A}.
\]
证明 $\scr{G}^p:=\oplus_{i=p}^{r}\Sym^i\scr{F}'|_{\Spec A}\otimes_{\OO_{\Spec A}}\Sym^{r-i}\scr{F}''|_{\Spec A}$ 给出一个\trWellDefined{}的（\trLocallyFree{}的）\trSubsheaf，它与所做的一切选取——例如诸基 $e_1,\dots,e_p$ 与 $f_1,\dots,f_q$，以及诸提升 $f_1',\dots,f_q'$——无关。）

\textbf{14.2.N.（有用的练习）}\quad 设 $0\to\scr{F}'\to\scr{F}\to\scr{F}''\to 0$ 是\trLocallyFreeSheaf{}的一个正合列。证明：对任意 $r$，$\wedge^r\scr{F}$ 有一个\trFiltration
\[
\wedge^r\scr{F}=\scr{G}^0\supset\scr{G}^1\supset\cdots\supset\scr{G}^r\supset\scr{G}^{r+1}=0,
\]
其\trSubquotient{}为
\[
\scr{G}^p/\scr{G}^{p+1}\cong(\wedge^p\scr{F}')\otimes(\wedge^{r-p}\scr{F}'')
\]
对每个 $p$ 成立。特别地，若这些层具有有限\trRank，则 $\det\scr{F}=(\det\scr{F}')\otimes(\det\scr{F}'')$。

\textbf{14.2.O.（练习）}\quad 设 $\scr{F}$ 是\trRank{}为 $n$ 的\trLocallyFreeSheaf。描述一个映射 $\wedge^r\scr{F}\times\wedge^{n-r}\scr{F}\to\wedge^n\scr{F}$，使之诱导一个\trIsomorphism{} $\wedge^r\scr{F}\tilde{\longrightarrow}(\wedge^{n-r}\scr{F})^{\vee}\otimes\wedge^n\scr{F}$。这称为一个\textbf{\trVectorBundle{}的\trPerfectPairing}（perfect pairing of vector bundles）。（若你了解\trVectorSpace{}的\trPerfectPairing，你看得出这为什么是它的一个推广吗？）日后在证明\trAlgClosedField{}上正则簇的\trHodge{}数对偶性时（\mbox{练习~21.5.L}），你或许会用上它。

\textbf{14.2.P.（练习：\trDeterminant{}\trLineBundle{}在正合列中表现良好）}\quad 设 $0\to\scr{F}_1\to\cdots\to\scr{F}_n\to 0$ 是 $X$ 上有限\trRank{}的\trLocallyFreeSheaf{}的一个正合列。证明：「行列式丛的交错积是平凡的」：
\[
\det(\scr{F}_1)\otimes\det(\scr{F}_2)^{\vee}\otimes\det(\scr{F}_3)\otimes\det(\scr{F}_4)^{\vee}\otimes\cdots\otimes\det(\scr{F}_n)^{(-1)^n}\cong\OO_X.
\]
（提示：把这个正合列打断成若干个\trShortExactSequence。用\mbox{练习~14.2.H(a)}证明：它们是\emph{有限\trRank{}的\trLocallyFreeSheaf}的\trShortExactSequence。然后用\mbox{练习~14.2.N}。）

\textbf{14.2.Q.（练习）}\quad 设 $\scr{F}$ 是一个\trQCSheaf。定义\trQCSheaf{} $\mathbf{T}^n\scr{F}$、$\Sym^n\scr{F}$ 与 $\wedge^n\scr{F}$。（一种可能：在每个仿射开集上描述它们，并利用重要练习~6.2.D 的刻画。）
